% !TEX root = ../../hw04.tex
% Author's handwritten solution: images/HW05/IMG_9276--IMG_9278.HEIC.
% Retain the tail-bound and monotone-partial-sum methods. The alternative
% direct-comparison argument on IMG_9278 is recorded in the final paragraph.
% Finite prefixes and partial sums, not unproved infinite sums, are compared.
% The author approved the editorial changes; review colouring is removed.
\begin{proof}[Solution]
  Recall the tail infima and suprema of the ratios.
  \[
    \alpha_n=\inf\left\{\frac{a_j}{b_j}\mid j\ge n\right\},
    \qquad
    \beta_n=\sup\left\{\frac{a_j}{b_j}\mid j\ge n\right\}.
  \]
  \begin{align*}
    r&=\sup_n\alpha_n=\lim_{n\to\infty}\alpha_n,\\
    R&=\inf_n\beta_n=\lim_{n\to\infty}\beta_n.
  \end{align*}
  We can think of $r$ and $R$ as long-term lower and upper bounds
  with the following precise meaning. A bound strictly
  below $r$ or strictly above $R$ holds on a sufficiently late tail.

  \emph{Part (i).} Since $r>0$, there is $c>0$
  with $c<r$. Choose a finite such $c$, also when
  $r=+\infty$. Since $r=\sup_n\alpha_n$, some $N\in\N$ has
  $\alpha_N>c$. Thus
  \[
    \forall n\ge N\quad \frac{a_n}{b_n}\ge c.
  \]
  Since $b_n>0$, this gives $a_n\ge cb_n$ for every $n\ge N$.

  Sum over this tail, where the comparison is known to hold. For $m\ge N$,
  \[
    \sum_{n=N}^{m}a_n\ge c\sum_{n=N}^{m}b_n.
  \]
  Removing the first $N-1$ terms does not change divergence.

  Since $b_n>0$, its partial sums are increasing and tend to $+\infty$.
  Therefore the partial sums of $\sum a_n$ also tend to $+\infty$.
  Hence $\sum_{n=1}^{\infty}a_n$ diverges.

  \emph{Part (ii).} Since $R<\infty$, choose $M>R$.
  The ratios are nonnegative, so $R\ge0$ and $M>0$.
  Since $R=\inf_n\beta_n$, some $N\in\N$ has $\beta_N<M$.
  It follows that
  \[
    \forall n\ge N\quad \frac{a_n}{b_n}\le M,
    \qquad 0\le a_n\le Mb_n.
  \]
  We know that $\sum_{n=1}^{\infty}b_n=L$ is finite.

  Put $A_m=\sum_{n=1}^{m}a_n$. For every $m\ge N$,
  \begin{align*}
    A_m
      &=\sum_{n=1}^{N-1}a_n+\sum_{n=N}^{m}a_n\\
      &\le\sum_{n=1}^{N-1}a_n+M\sum_{n=N}^{m}b_n\\
      &\le\sum_{n=1}^{N-1}a_n+ML.
  \end{align*}
  The finite prefix is zero when $N=1$. Earlier partial sums are
  bounded by the same constant because the terms are nonnegative.

  This establishes a finite upper bound. Since
  $a_n\ge0$, the partial sums are monotone nondecreasing.
  Therefore, by the monotone convergence theorem,
  $A_m$ tends to a finite limit.
  Hence $\sum_{n=1}^{\infty}a_n$ converges.

  Direct comparison gives the same conclusion
  on the tail. Adding the finite prefix preserves
  convergence.
\end{proof}
